% Written by unit TR-13.
\section{Notation}
\label{sec:notation}

The method sections were written one family of models at a time, and a symbol that is natural in one family is used with another meaning in the next:
the letter $\alpha$ is the tail probability of a value at risk in the market section and the level of a test in the validation sections. This section fixes the meaning that a symbol keeps across the report, and the table lists the symbols that two or more sections share. A section that
needs a symbol in another sense, or a symbol that only it uses, defines it where it first appears, and the definition there prevails; the symbols that carry a
different meaning in one section are marked in the last column of \cref{tab:notation}.

Four conventions apply to every section. First, probabilities and rates are decimals in the equations and percentages in the text and the tables, with the unit named; a
rate is a fraction of the exposure it applies to, and a loss is an amount in the currency of the case (United States dollars, euro or Chilean peso, with the unit of thousands or millions
named in the caption). Second, a hat marks an estimate and a star marks a truth that is known because the data were generated, so that $\hat\theta-\theta^\ast$ is an error that can be
measured. Third, the index $i$ runs over obligors, loans or models, the index $t$ over days, years or periods, the index $j$ over grades, buckets or segments and the index $b$ over
the repetitions of a simulation. Fourth, a statistic that is compared with a threshold is written $T$, its p-value $p$, and the light assigned to it follows
the rule of \cref{eq:vt-light}; the levels $\alpha_{\mathrm{amber}}$ and $\alpha_{\mathrm{red}}$ belong to the policy file of the engine and are conventions.

The default probability, the loss given default and the exposure at default are the three parameters of the credit sections, and their product is the expected loss that the
provisions and the capital sections build on. The standard normal distribution function $\Phi$ and its inverse enter through the one-factor model of Vasicek, which links the
ratings, provisions, capital, counterparty and validation sections: an obligor defaults when its asset value falls to $\Phi^{-1}(\PD)$, and the asset value mixes the systematic
factor $Z$ and an idiosyncratic factor with the weights set by the asset correlation $\rho$. The market, counterparty, banking-book and operational sections add the risk
measures (value at risk and expected shortfall), the exposure profile and its summaries, the changes in economic value and net interest income, and the tail index of the severity.
The last group of symbols belongs to the sequential monitors of the proposed methods, where an e-process $E_t$ turns the evidence of each day into a bet whose product is a
valid test at any stopping time.

The units follow the International System, with the financial units named in each case. Quantities in the tables are rounded for reading; the macros that the text
reads carry the values of the committed artifacts to the digits the artifact stores, and no rounded value is fed back into a computation.

\begin{table}[htbp]
\centering
\footnotesize
\setlength{\tabcolsep}{4pt}
\begin{tabularx}{\textwidth}{>{\raggedright\arraybackslash}p{2.6cm}>{\raggedright\arraybackslash}X>{\raggedright\arraybackslash}p{2.7cm}>{\raggedright\arraybackslash}p{2.9cm}}
\toprule
Symbol & Meaning & Unit & Main sections \\
\midrule
$\PD,\ \LGD,\ \EAD$ & Probability of default over a stated horizon, loss given default as a share of the exposure, and exposure at default & decimal, decimal, currency & \cref{sec:ratings,sec:provisions,sec:lgd-ead,sec:credit-capital} \\
$\ECL$ & Expected credit loss of a loan or a book under the staging and the horizon of the accounting standard & currency & \cref{sec:provisions} \\
$\Phi,\ \varphi$ & Standard normal distribution function and density & dimensionless & \cref{sec:ratings,sec:credit-capital,sec:market,sec:ccr} \\
$Z,\ \varepsilon_i$ & Systematic factor and idiosyncratic factor of obligor $i$ in the one-factor model, both standard normal and independent & dimensionless & \cref{sec:ratings,sec:credit-capital} \\
$\rho$ & Asset correlation of the one-factor model & decimal in $(0,1)$ & \cref{sec:ratings,sec:credit-capital,sec:validation-tests} \\
$N,\ D$ & Obligors at the start of a year and defaults during it (a different meaning in \cref{sec:ccr}, where $N$ is the normal distribution function) & count & \cref{sec:ratings,sec:validation-tests} \\
$q$ & Tail probability or confidence level of a risk measure & decimal & \cref{sec:credit-capital,sec:market,sec:aggregation-stress} \\
$\VaR_q,\ \ES_q$ & Value at risk and expected shortfall at level $q$ of a loss distribution or a return series & currency or decimal return & \cref{sec:credit-capital,sec:market,sec:aggregation-stress} \\
$\mathrm{AUC}$ & Area under the receiver operating characteristic curve, the probability that a defaulter is ranked above a non-defaulter & decimal & \cref{sec:credit-scoring,sec:validation-tests} \\
$\mathrm{EE},\ \mathrm{EEPE},\ \mathrm{CVA}$ & Expected exposure profile, effective expected positive exposure, and credit valuation adjustment of a netting set & currency & \cref{sec:ccr} \\
$\mathrm{EVE},\ \mathrm{NII}$ & Economic value of equity and net interest income of the banking book, with $\Delta$ for the change under a shock & currency & \cref{sec:alm} \\
$\xi$ & Tail index of a heavy-tailed severity (the shape of the generalised Pareto family); a value of 1 or more implies an infinite mean & dimensionless & \cref{sec:oprisk,sec:market} \\
$T,\ p,\ H_0$ & Statistic of a validation test, its p-value and its null hypothesis & as the test defines & \cref{sec:validation-tests,sec:problem} \\
$\alpha_{\mathrm{amber}},\ \alpha_{\mathrm{red}}$ & Levels of the policy that colour a p-value test & decimal & \cref{sec:validation-tests} \\
$B$ & Number of independent repetitions of a simulation, each with its own random stream & count & \cref{sec:problem,sec:protocol,sec:validation-tests} \\
$E_t$ & E-process, a non-negative process with unit start whose expectation under the null is at most one at every stopping time & dimensionless & \cref{sec:proposed} \\
$\E,\ \Prob,\ \Var,\ \indic$ & Expectation, probability, variance and indicator function & as the argument & all \\
$\hat\theta,\ \theta^\ast$ & Estimate of a parameter and its true value on a generator of known truth & as the parameter & all \\
\bottomrule
\end{tabularx}
\caption{Symbols that two or more sections of the report share, with the meaning they keep, their unit and the sections that use them; a symbol marked in the meaning column with a different use in another section is redefined there, and the definition in that section prevails. A symbol that only one section uses is defined where that section introduces it and is not listed here.}
\label{tab:notation}
\end{table}

\Cref{tab:notation} does not list the symbols of a single method, such as the parameters of the Hull--White model of the counterparty section or the shape functions of the explainable boosting
machine; each is defined at its first use, in the section that needs it.
